\subsection*{用 \texttt{\#} 合成两个有界变量以分块}

第 3.10.2 节 RGB 转灰度示例将工作实例索引 q 与顺序图像索引 n 组合起来。\co@q # n@ 为每个实例分配连续图像。N = 18、共有六个实例时，内层上界是 3，所以实例 1 处理图像 3、4、5。下面沿用这一划分。

共有六个工作实例时，q 的上界是 6，n 的上界是 N / 6。组合后的上界为 N，值为 \co@q * (N / 6) + n@。其余索引 h、w 选择空间分块，通道维保持完整。因此，实例分配与局部遍历共同确定源分块及对应输出位置。

交换操作数会改变分配方式。N = 18 时，随着 n 递增，\co@n # q@ 会让实例 1 处理图像 1、7、13，而 \co@q # n@ 选择的是连续的图像 3、4、5。运算符将这种调度选择直接写入数据表达式，实际性能取决于所得访问在目标平台上的表现。
